% GENERATED FILE. Edit canonical lessons, then run npm run generate:books.
% canonical-source-hash: fnv1a64:d2f7ca9ec5c10896
% canonical-lessons: JA-C516-shusse-suru, JA-C516-norikoeru, JA-C516-zenryoku, JA-C516-yaruki, JA-C516-idomu

\chapter{Getting Ahead}
\label{ch:ja-getting-ahead}
\hlchaptermodality{516}

\begin{chapteropening}
\textbf{By the end of this chapter:} \emph{I can talk about success, failure and getting through.}

Complete the last lesson of chapter 516: I can talk about success, failure and getting through.
\end{chapteropening}

\section[\texorpdfstring{\ja{しゅっせする} (shusse suru)}{しゅっせする (shusse suru)}]{\ja{しゅっせする} (shusse suru) — to get ahead (in life)}

\label{lesson:JA-C516-shusse-suru}



\begin{quote}
Before the new one: say the Japanese for grades, then the Japanese for a cram school (for entrance exams).
\end{quote}

\subsection*{You'll want to know: \ja{しゅっせする}}
\textbf{\ja{しゅっせする}} — \emph{shusse suru} — \textquotedblleft{}to get ahead (in life)\textquotedblright{}.

\begin{quote}
\textbf{\ja{もし} \ja{かれが} \ja{しゅっせしたら、しゃちょうに} \ja{なるでしょう。}} — \emph{moshi kare ga shusseshitara, shachō ni naru deshō} — \textquotedblleft{}If he gets ahead, he will probably become company president.\textquotedblright{}
\end{quote}

Polite \textbf{\ja{しゅっせします}}, \emph{shusseshimasu}; a \ja{する} verb. Past (polite): \textbf{\ja{あにが} \ja{かいしゃで} \ja{しゅっせしました。}}, \emph{ani ga kaisha de shusseshimashita}.

\begin{cousinweb}
\textbf{\ja{しゅっせする}} (\emph{shusse suru}) is the Sino-Japanese (kango) noun \emph{shusse} with the native verb \emph{suru}, \textquotedblleft{}to do\textquotedblright{}: two Chinese characters read with their on'yomi, readings borrowed from Chinese. They are \emph{shu}, \textquotedblleft{}exit\textquotedblright{} (Middle Chinese \emph{tsyhwit}) and \emph{se}, \textquotedblleft{}generation\textquotedblright{} (\emph{syejH}) (KANJIDIC2; Middle Chinese in Baxter's transcription).
\end{cousinweb}

\begin{grammarlens}[title={What you've built}]
One more word for saying what happens if.
\end{grammarlens}

\subsection*{Your turn}
\begin{itemize}
\raggedright
  \item \emph{Say it:} \emph{shusse suru}
  \item \emph{Say it:} \emph{shusse suru}, once more
  \item \emph{Say it:} \emph{moshi kare ga shusseshitara, shachō ni naru deshō} — If he gets ahead, he will probably become company president.
  \item \emph{Recall:} say the Japanese for grades, then the Japanese for a cram school (for entrance exams), then say \emph{shusse suru} again
\end{itemize}

\begin{tcolorbox}[breakable,colback=teal!4,colframe=teal!35!black,title={Before you move on}]
What does \ja{しゅっせする} mean? (\textquotedblleft{}To get ahead (in life)\textquotedblright{}.) Say it once more.
\end{tcolorbox}
\section[\texorpdfstring{\ja{のりこえる} (norikoeru)}{のりこえる (norikoeru)}]{\ja{のりこえる} (norikoeru) — to get over; to overcome}

\label{lesson:JA-C516-norikoeru}



\begin{quote}
Before the new one: say the Japanese for a cram school (for entrance exams), then the Japanese for to get ahead (in life).
\end{quote}

\subsection*{You'll want to know: \ja{のりこえる}}
\textbf{\ja{のりこえる}} — \emph{norikoeru} — \textquotedblleft{}to get over; to overcome\textquotedblright{}.

\begin{quote}
\textbf{\ja{もし} \ja{この} \ja{しけんを} \ja{のりこえたら、そつぎょうです。}} — \emph{moshi kono shiken o norikoetara, sotsugyō desu} — \textquotedblleft{}If I get through this exam, I graduate.\textquotedblright{}
\end{quote}

Polite \textbf{\ja{のりこえます}}, \emph{norikoemasu}; an ichidan verb. Past (polite): \textbf{\ja{かれが} \ja{びょうきを} \ja{のりこえました。}}, \emph{kare ga byōki o norikoemashita}.

\begin{cousinweb}
\textbf{\ja{のりこえる}} (\emph{norikoeru}) is a native Japanese word (wago), inherited rather than borrowed from Chinese, built from \emph{noru}, \textquotedblleft{}to ride\textquotedblright{}, and \emph{koeru}, \textquotedblleft{}to cross\textquotedblright{}; its kanji was attached later as a kun'yomi (JMdict, KANJIDIC2).
\end{cousinweb}

\begin{grammarlens}[title={What you've built}]
One more word for saying what happens if.
\end{grammarlens}

\subsection*{Your turn}
\begin{itemize}
\raggedright
  \item \emph{Say it:} \emph{norikoeru}
  \item \emph{Say it:} \emph{norikoeru}, once more
  \item \emph{Say it:} \emph{moshi kono shiken o norikoetara, sotsugyō desu} — If I get through this exam, I graduate.
  \item \emph{Recall:} say the Japanese for a cram school (for entrance exams), then the Japanese for to get ahead (in life), then say \emph{norikoeru} again
\end{itemize}

\begin{tcolorbox}[breakable,colback=teal!4,colframe=teal!35!black,title={Before you move on}]
What does \ja{のりこえる} mean? (\textquotedblleft{}To get over; to overcome\textquotedblright{}.) Say it once more.
\end{tcolorbox}
\section[\texorpdfstring{\ja{ぜんりょく} (zenryoku)}{ぜんりょく (zenryoku)}]{\ja{ぜんりょく} (zenryoku) — all one's strength; one's utmost}

\label{lesson:JA-C516-zenryoku}



\begin{quote}
Before the new one: say the Japanese for to get ahead (in life), then the Japanese for to get over.
\end{quote}

\subsection*{You'll want to know: \ja{ぜんりょく}}
\textbf{\ja{ぜんりょく}} — \emph{zenryoku} — \textquotedblleft{}all one's strength; one's utmost\textquotedblright{}.

\begin{quote}
\textbf{\ja{もし} \ja{ぜんりょくで} \ja{はしったら、かれは} \ja{まにあうでしょう。}} — \emph{moshi zenryoku de hashittara, kare wa maniau deshō} — \textquotedblleft{}If he runs with all his strength, he will make it.\textquotedblright{}
\end{quote}

\begin{cousinweb}
\textbf{\ja{ぜんりょく}} (\emph{zenryoku}) is Sino-Japanese (kango): two Chinese characters read with their on'yomi, readings borrowed from Chinese. They are \emph{zen}, \textquotedblleft{}whole\textquotedblright{} (Middle Chinese \emph{dzjwen}) and \emph{ryoku}, \textquotedblleft{}power\textquotedblright{} (\emph{lik}) (KANJIDIC2; Middle Chinese in Baxter's transcription).
\end{cousinweb}

\begin{grammarlens}[title={What you've built}]
One more word for saying what happens if.
\end{grammarlens}

\subsection*{Your turn}
\begin{itemize}
\raggedright
  \item \emph{Say it:} \emph{zenryoku}
  \item \emph{Say it:} \emph{zenryoku}, once more
  \item \emph{Say it:} \emph{moshi zenryoku de hashittara, kare wa maniau deshō} — If he runs with all his strength, he will make it.
  \item \emph{Recall:} say the Japanese for to get ahead (in life), then the Japanese for to get over, then say \emph{zenryoku} again
\end{itemize}

\begin{tcolorbox}[breakable,colback=teal!4,colframe=teal!35!black,title={Before you move on}]
What does \ja{ぜんりょく} mean? (\textquotedblleft{}All one's strength; one's utmost\textquotedblright{}.) Say it once more.
\end{tcolorbox}
\section[\texorpdfstring{\ja{やるき} (yaruki)}{やるき (yaruki)}]{\ja{やるき} (yaruki) — the will to do something; motivation}

\label{lesson:JA-C516-yaruki}



\begin{quote}
Before the new one: say the Japanese for to get over, then the Japanese for all one's strength.
\end{quote}

\subsection*{You'll want to know: \ja{やるき}}
\textbf{\ja{やるき}} — \emph{yaruki} — \textquotedblleft{}the will to do something; motivation\textquotedblright{}.

\begin{quote}
\textbf{\ja{もし} \ja{やるきが} \ja{あったら、かれは} \ja{しゅっせするでしょう。}} — \emph{moshi yaruki ga attara, kare wa shusse suru deshō} — \textquotedblleft{}If he had the drive, he would get ahead.\textquotedblright{}
\end{quote}

\begin{cousinweb}
\textbf{\ja{やるき}} (\emph{yaruki}) mixes the two layers, native and Sino-Japanese: native \emph{yaru}, \textquotedblleft{}to do\textquotedblright{}, and Sino-Japanese \emph{ki}, \textquotedblleft{}spirit\textquotedblright{} (Middle Chinese \emph{khj+jH}), borrowed from Chinese (JMdict, KANJIDIC2; Middle Chinese in Baxter's transcription).
\end{cousinweb}

\begin{grammarlens}[title={What you've built}]
One more word for saying what happens if.
\end{grammarlens}

\subsection*{Your turn}
\begin{itemize}
\raggedright
  \item \emph{Say it:} \emph{yaruki}
  \item \emph{Say it:} \emph{yaruki}, once more
  \item \emph{Say it:} \emph{moshi yaruki ga attara, kare wa shusse suru deshō} — If he had the drive, he would get ahead.
  \item \emph{Recall:} say the Japanese for to get over, then the Japanese for all one's strength, then say \emph{yaruki} again
\end{itemize}

\begin{tcolorbox}[breakable,colback=teal!4,colframe=teal!35!black,title={Before you move on}]
What does \ja{やるき} mean? (\textquotedblleft{}The will to do something; motivation\textquotedblright{}.) Say it once more.
\end{tcolorbox}
\section[\texorpdfstring{\ja{いどむ} (idomu)}{いどむ (idomu)}]{\ja{いどむ} (idomu) — to tackle (a challenge); to take on}

\label{lesson:JA-C516-idomu}



\begin{quote}
Before the new one: say the Japanese for all one's strength, then the Japanese for the will to do something.
\end{quote}

\subsection*{You'll want to know: \ja{いどむ}}
\textbf{\ja{いどむ}} — \emph{idomu} — \textquotedblleft{}to tackle (a challenge); to take on\textquotedblright{}.

\begin{quote}
\textbf{\ja{もし} \ja{まけたら、らいねん} \ja{もう} \ja{いちど} \ja{いどみます。}} — \emph{moshi maketara, rainen mō ichido idomimasu} — \textquotedblleft{}If I lose, I will take it on once more next year.\textquotedblright{}
\end{quote}

Polite \textbf{\ja{いどみます}}, \emph{idomimasu}; a godan verb. Past (polite): \textbf{\ja{かれが} \ja{たかい} \ja{やまに} \ja{いどみました。}}, \emph{kare ga takai yama ni idomimashita}.

\begin{cousinweb}
\textbf{\ja{いどむ}} (\emph{idomu}) is a native Japanese word (wago), inherited rather than borrowed from Chinese; its kanji was attached later as a kun'yomi, and its older origin is not settled (JMdict, KANJIDIC2).
\end{cousinweb}

\begin{grammarlens}[title={What you've built}]
One more word for saying what happens if.
\end{grammarlens}

\subsection*{Your turn}
\begin{itemize}
\raggedright
  \item \emph{Say it:} \emph{idomu}
  \item \emph{Say it:} \emph{idomu}, once more
  \item \emph{Say it:} \emph{moshi maketara, rainen mō ichido idomimasu} — If I lose, I will take it on once more next year.
  \item \emph{Recall:} say the Japanese for all one's strength, then the Japanese for the will to do something, then say \emph{idomu} again
\end{itemize}

\begin{tcolorbox}[breakable,colback=teal!4,colframe=teal!35!black,title={Before you move on}]
What does \ja{いどむ} mean? (\textquotedblleft{}To tackle (a challenge); to take on\textquotedblright{}.) Say it once more.
\end{tcolorbox}
